\clearpage\fancyhead[L]{\small PALLAS.}\SourceLeaf{381}
\begin{center}\textbf{PALLAS.}\par\vspace{4pt}\hrule\vspace{4pt}\fontsize{10}{12}\selectfont Astronomisches Jahrbuch für 1817. Seite 212\;..\;215. Berlin 1814.\par\vspace{4pt}\hrule\vspace{4pt}\end{center}

\begin{flushright}
Göttingen den 29. Mai 1814.
\end{flushright}

\par\indent Ich habe Ihnen Verehrtester Freund noch meinen verbindlichsten Dank nachzuholen für das schätzbare Geschenk Ihres astronomischen Jahrbuchs für 1816; ich habe denselben aufgeschoben, weil ich ihn mit einigen Beiträgen für das neue Jahrbuch zu begleiten wünschte; indessen ist die Beendigung dessen, was ich Ihnen über die letzte Vestaopposition schicken wollte, bisher noch aufgehalten, ich werde dies aber künftig noch nachholen.

\Needspace{14\baselineskip}\par\indent Ueber die letzte Junoopposition habe ich Ihnen zwar schon einiges gemeldet, da ich aber jetzt noch verschiedenes darüber beifügen kann, und nicht genau mehr weiss, wie viel ich Ihnen schon geschrieben habe, so will ich hier alles Wesentliche zusammen stellen. Das Wetter erlaubte mir nur eine Beobachtung um die Zeit der Opposition, die ich indess für sehr gut halte; erst beinahe drei Wochen nachher konnte ich zur Controlle noch eine zweite machen, welche letztere ich nach Herrn \textsc{Nicolai}'s Reduction beifüge:

\begin{center}\fontsize{10}{12}\selectfont\setlength{\tabcolsep}{3pt}\renewcommand{\arraystretch}{1.15}\begin{tabular}{l|c|rrr|rrr}
1813 & Mittlere Zeit & \multicolumn{3}{c|}{Scheinbare AR.} & \multicolumn{3}{c}{Südl. Abweich.}\\\hline
Nov. 19 & $9^{\mathrm u}\;46^{\mathrm m}\;46^{\mathrm s}$ & $60^\circ$ & $30'$ & $35''0$ & $3^\circ$ & $5'$ & $9''3$\\
Dec. 9 & $7\;9\;38$ & 56 & 56 & 47.4 & 3 & 56 & 35.3\\
\end{tabular}\end{center}

\Needspace{14\baselineskip}\par\indent Ich hatte zwar schon auf meine erste Beobachtung die vorläufige Berechnung der Opposition gegründet, allein da theils meine Rechnung nur flüchtig gemacht war, theils mir nachher noch von \textsc{Bessel} seine Meridianbeobachtungen mitgetheilt wurden (welche Sie ohne Zweifel von ihm selbst erhalten haben werden), so veranlasste ich einen geschickten und talentvollen Zuhörer von mir, Hrn. \textsc{Möbius}, die Resultate hiernach neu zu berichtigen. Er findet die Opposition:

\begin{center}1813 November 19. $18^{\mathrm u}\;12^{\mathrm m}\;0^{\mathrm s}$ Mittlere Zeit in Göttingen\end{center}
\begin{center}\fontsize{10}{12}\selectfont\setlength{\tabcolsep}{2pt}\renewcommand{\arraystretch}{1.15}\begin{tabular}{lcr}
Wahre Länge & $.\;.\;.\;.\;.\;.\;.\;.\;.\;.\;.\;.\;.\;.$ & $57^\circ\;34'\;2''1$\\
Geocentrische Breite & $.\;.\;.\;.\;.\;.\;.\;.\;.\;.\;.$ & $23\;18\;46.1\;\mathrm{S.}$\\
\end{tabular}\end{center}

\Needspace{14\baselineskip}\par\indent Er hat zugleich dieser und den drei vorhergehenden Oppositionen von 1810, 1811, 1812 die Elemente angepasst, und folgendes Resultat erhalten:

\begin{center}\fontsize{9}{11}\selectfont\setlength{\tabcolsep}{2pt}\renewcommand{\arraystretch}{1.15}\begin{tabular}{lcr}
Epoche 1810 Meridian von Göttingen & $.\;.\;.\;.\;.\;.\;.$ & $95^\circ\;29'\;53''2$\\
Tägliche mittlere tropische Bewegung & $.\;.\;.\;.\;.\;.\;.$ & $812''7140$\\
Länge der Sonnennähe (1810 sid. ruh.) & $.\;.\;.\;.\;.\;.\;.$ & $53^\circ\;6'\;43''0$\\
Länge des aufsteigenden Knoten (1810) & $.\;.\;.\;.\;.\;.$ & $171\;6\;45.0$\\
Excentricitätswinkel & $.\;.\;.\;.\;.\;.\;.\;.\;.\;.\;.\;.\;.$ & $14\;43\;9.5$\\
Neigung der Bahn & $.\;.\;.\;.\;.\;.\;.\;.\;.\;.\;.\;.\;.\;.$ & $13\;4\;12.9$\\
Logarithm. der halben grossen Axe & $.\;.\;.\;.\;.\;.\;.$ & $0.4267616$\\
\end{tabular}\end{center}

\par\indent Ueberdies hat dieser geschickte Rechner für die nächste Erscheinung der Juno eine Ephemeride berechnet, wovon ich das Vergnügen habe, Ihnen eine Abschrift beizulegen.

\SourceLeaf{382}\par\indent Für die Pallas hat Hr. \textsc{Nicolai} (welcher seit einem halben Jahre als Gehülfe bei der
Seeberger Sternwarte angestellt ist) gleichfalls eine Ephemeride berechnet, und mir beiliegende
Copie davon für Ihr Jahrbuch geschickt. Zum Auffinden des Planeten (der ohnehin diesmal ziemlich
hell sein wird) ist dieselbe gewiss überflüssig genau, ob sie gleich nach den frühern Elementen
berechnet ist, ohne deren Aenderung durch die fortwirkenden Störungen Jupiters zu berücksichtigen.
Inzwischen ist es doch interessant, den Einfluss welchen diese haben werden \emph{im Voraus}
anzugeben, und ich habe daher Hrn. \textsc{Möbius} diese ziemlich mühsame Arbeit aufgetragen. Das
Resultat ist, dass um die Zeit der Opposition die gerade Aufsteigung in Hrn. \textsc{Nicolai}'s
Ephemeride um 1$\frac{1}{2}$ Minuten vermindert werden muss, während die Declination nur eine
Verminderung von wenigen Secunden erleidet. Für die Opposition selbst erhält er im Voraus
folgendes:

\begin{center}1814 October 25. $12^{\mathrm u}\;39^{\mathrm m}\;50^{\mathrm s}$ Göttinger Zeit.\end{center}
\begin{center}\fontsize{10}{12}\selectfont\setlength{\tabcolsep}{2pt}\renewcommand{\arraystretch}{1.15}\begin{tabular}{lcr}
Wahre Länge & $.\;.\;.\;.\;.\;.\;.\;.\;.\;.$ & $31^\circ\;58'\;28''4$\\
Heliocentrische Breite & $.\;.\;.\;.\;.\;.\;.$ & $23\;38\;32.7\;\mathrm{S.}$\\
\end{tabular}\end{center}

\noindent (die geocentrische hat er vergessen mir mitzutheilen, und ich kann sie daher, da er sich jetzt
nicht mehr hier aufhält, zur Vergleichung mit Hrn. \textsc{Nicolai}'s Resultat nicht beifügen.)

\par\indent Ich habe jetzt angefangen, auch die Störungen der Pallas durch Mars zu berechnen. Sie werden
zwar fast alle sehr klein, allein doch zahlreich. Alle Gleichungen durch Jupiter, Saturn und
Mars, die über $0''1$ gehen, werden eine Anzahl von mehr als 1000 ausmachen.

\Needspace{10\baselineskip}\par\indent Vor ein Paar Tagen habe ich aus München einen trefflichen Heliometer von 43 Zoll Brennweite,
34 Linien Oeffnung erhalten. Der geniale Künstler hat unter andern die wichtige Verbesserung
angebracht, dass er ihn zum Repetiren eingerichtet hat, indem er beide Objectivhälften
unabhängig von einander beweglich gemacht hat. Man vermeidet dadurch ganz den nachtheiligen
Einfluss, welchen kleine locale Ungleichheiten der Schraube hervorbringen. Die Präcision des
Fernrohrs, ist, wie man dies von \textsc{Frauenhofer}'s Arbeiten gewohnt ist, ausserordentlich
gross.
